\section{Related Work}
\label{sec:related}

\begin{intuition}
This paper sits at the intersection of four literatures: (i) retrieval-
augmented generation and vector-database compression; (ii) the
cognitive-science theory of memory consolidation; (iii) the geometry
of high-dimensional embedding spaces; and (iv) continual learning.
Each has produced a piece of the consolidation problem, and this
section places our contribution against them.
\end{intuition}

\subsection{Retrieval-augmented generation and vector compression}
\label{sec:related:rag}

Retrieval-augmented generation (RAG)~\cite{lewisRAG2020,
guu2020realm, karpukhin2020dpr, borgeaud2022retro, izacardAtlas2023,
izacard2021fid} adds a vector-database retrieval step to a
pretrained language model: given a query, fetch the top-$k$ relevant
passages and condition the reader on their concatenation. The
dominant engineering concern in deployed RAG systems is the
\emph{size} of the vector index. A billion-passage corpus encoded at
$768$--$1536$ dimensions is multiple terabytes in raw form; every
production system therefore applies some lossy compression.

The classical compression toolkit splits into two families. The
\emph{learned-quantization} family, dominated by Product
Quantization~\cite{jegouPQ2011} and its Optimized PQ
variant~\cite{ge2013optimized}, stores each vector as a short string of
codebook indices, typically at $4$--$32$ bytes per vector. PQ is
implemented in FAISS~\cite{johnson2019faiss} and is the default
compression for billion-scale retrieval. Simpler baselines include
locality-sensitive hashing~\cite{charikarLSH2002, datarLSH2004},
dimensionality reduction, and the HNSW graph pruning technique
built into most vector databases~\cite{malkovHNSW2018}. The
\emph{consolidation} family replaces a cluster of similar vectors
with one or a few representatives: centroid summarisation, medoid
selection, or selective pruning.

The two families have not, to our knowledge, been compared on a
common Pareto frontier before this work. The consolidation family
has been studied primarily in cognitive-science contexts (see
below), while the compression family has been studied in
engineering contexts. E3 (\S\ref{sec:e3}) is the first apples-to-apples
comparison, showing that consolidation Pareto-dominates learned
quantization in the low-$\deff$ corner (five of six corpora) and
the frontiers cross only at high $\deff$.

\subsection{Cognitive consolidation}
\label{sec:related:cls}

The idea that memory systems replace dense experience traces with
sparse representatives runs through half a century of cognitive
neuroscience. The standard-model framing is \emph{Complementary
Learning Systems}~\cite{mcclellandConsolidation1995,
alvarezSquireConsolidation1994, oreillyHippocampus2014, marrSimple1971},
which posits that the hippocampus stores episodic traces quickly and
noisily, while the neocortex slowly abstracts those traces into
semantic representatives through offline replay~\cite{wilson1994reactivation,
robinsReplay1995}. The CLS framework motivates a vast computational
literature on consolidation algorithms --- most of which treat
consolidation as parameter-update rather than index-update
\cite{mccloskey1989cf, kirkpatrickEWC2017, chaudhryGEM2019,
tulvingEpisodic1972}. Our contribution is to show that the
\emph{vector-index} version of consolidation, which is the form used
in production RAG, obeys the same capacity-bound structure as the
parameter-update version and can be analysed with the same geometric
tools.

\subsection{Geometry of embedding spaces and effective rank}
\label{sec:related:geometry}

The effective rank of a matrix (or a set of vectors) is an
entropy-weighted measure of its spectral spread, formalised by Roy
and Vetterli~\cite{royVetterli2007}. In high dimensions, the
cap-volume concentration phenomena~\cite{ball1997elementary,
milman1986asymptotic, ledoux2001concentration, vershynin2018hdp,
beyer1999nn} constrain how much of the sphere a small neighbourhood
can occupy; the Berry-Esseen machinery~\cite{esseen1942} then gives
uniform rates for the Gaussian approximation that underlies our
bound. The first paper in the trilogy~\cite{geometryOfForgetting2026}
measured effective dimensions of production encoders (SBERT,
MiniLM, MPNet, BGE-base/large, E5-large, Nomic, SimCSE) and found
$\deff \approx 16$ across the board; the second paper raised this
into a lower bound on identity retrieval under noise. The present
paper extends the same geometric framework to the compressed case,
where the perturbation is not noise but a consolidator.

Production sentence encoders~\cite{reimers2019sbert, gao2021simcse,
wang2022e5, wang2020minilm, song2020mpnet, xiao2023bge, nussbaum2024nomic}
are the empirical objects on which the theorem is evaluated. Each
has its own training recipe (contrastive, self-supervised,
retrieval-fine-tuned) but the effective-dimensionality and
cap-volume phenomena are nearly encoder-independent across the
six encoders we test (\S\ref{sec:template_collapse}). This is consistent with
the hypothesis that the geometric structure is a property of the
\emph{semantic space}, not the individual encoder.

\subsection{Continual learning and catastrophic forgetting}
\label{sec:related:cl}

Catastrophic forgetting~\cite{mccloskey1989cf} is the well-known
failure mode of sequential neural-network training: learning a new
task overwrites the parameters needed for an old one. The
continual-learning literature~\cite{kirkpatrickEWC2017,
chaudhryGEM2019, robinsReplay1995} responds with replay buffers
(revisit old data), regularisation (EWC, MAS), or architectural
separation. The parallel between consolidation in RAG and replay in
continual learning is the direct inspiration for the term
``consolidation'' in our setting: in both cases, the operator reduces
a dense record of experience to a sparser one in the interest of
memory cost. The cap-volume bound we prove should be viewed as a
\emph{necessary condition} that any replay compression scheme must
satisfy; E5's $1$M scaling result confirms that the bound is not a
small-$n$ artefact.

\subsection{QA benchmarks and downstream evaluation}
\label{sec:related:qa}

The downstream E7 evaluation uses three QA benchmarks with
distinctive cluster structures: Natural Questions~\cite{kwiatkowski2019nq},
HotpotQA~\cite{yang2018hotpotqa}, and PopQA~\cite{mallen2023popqa};
we also test MS MARCO passage ranking~\cite{nguyen2016msmarco} in
E2/E4. The SQuAD v$1.1$ text-normalisation protocol
~\cite{rajpurkar2016squad} is used for the EM/F$1$ scoring.
Llama-$3.1$-$70$B-Instruct~\cite{llama31} is served through
vLLM~\cite{kwon2023vllm} as the reader. The DRM template
corpus~\cite{roediger1995drm} is a synthetic paraphrase benchmark
adapted from the original Deese-Roediger-McDermott false-memory
paradigm; it gives us a controlled high-$\theta$ regime for
operator isolation.

\section{Discussion}
\label{sec:discussion}

\begin{intuition}
\textbf{The shape of the result.} We have formulated a mathematically
explicit version of a simple question --- when can a cluster of
embeddings be replaced by a sparse set of representatives without
losing its identity? --- and used it to organise a body of
experimental evidence that previously looked like a loose collection
of empirical findings about RAG compression. The theorem is a
regime split: $\dbar < \thetap$ (tight) versus $\dbar \geq \thetap$
(spread). The unexpected empirical consequence of the split is that
on every real-text corpus we tested, the simplest one-vector
summary --- the cluster centroid --- dominates everything we or
anyone else might build on top of it. Our own adaptive scheme, GAC,
is dominated by the fixed centroid on five of six corpora, and even
the in-hindsight oracle cannot meaningfully beat it. We read this
as confirmation that real English text at retrieval thresholds near
$0.8$ sits inside the tight regime the law describes.
\end{intuition}

\subsection{The law, stated in plain language}
\label{sec:disc:law}

Restated without symbols: an embedding cluster can be safely
replaced by a one-vector summary exactly when the cluster's mean
cosine spread is smaller than the retrieval threshold's orthogonal
complement, $1 - \theta$. When that condition holds, the cap around
the centroid contains every cluster member and retrieval is
recovered by the summary at negligible byte cost. When it fails, no
single vector suffices; in fact, no cluster-aware summary at all
can recover what the full cluster would have provided, and learned
vector quantization (which works per-vector rather than per-cluster)
overtakes consolidation entirely.

This is a statement we believe to be the first mathematically
explicit formulation of the consolidation-compressibility law. It is
not the final form. The bound is isotropic and becomes vacuous in
the extreme-coherence corner of the spread regime
(\S\ref{sec:e1:calibration}); we see that as an invitation for
future work rather than a flaw in the framing.

\subsection{The centroid lesson}
\label{sec:disc:centroid}

On every real-text corpus we tested, the centroid and its
importance-weighted twin dominate every adaptive scheme we could
build, including the one we built specifically to beat it (GAC).
The oracle ceiling on any router is within a few EM points of the
fixed centroid, and on Wikipedia sections and NQ the oracle is
\emph{worse} than the fixed centroid. We read this as evidence that
real English text at $\theta = 0.8$ sits well inside the tight
regime of the law, and that the law's headline prediction in that
regime is that the centroid is already near-optimal. The GAC router
is therefore best understood as the instrument that made the
measurement, not as the method that the paper advocates.

The centroid-dominance result also has a deployment implication.
Most RAG systems we are aware of implement some variant of an
adaptive compression policy. If the corpus sits in the tight regime
(readily checkable from $\deff_{\mathrm{local}}$ and $\dbar$), no
adaptive policy will recover meaningful accuracy from a fixed
centroid. We suggest running the cheap geometric check before
investing in a router.

\subsection{Scope of our claims}
\label{sec:disc:scope}

We take care to state the scope narrowly. The theorem and every
experimental claim in this paper apply to \emph{embedding
consolidation in contrastive unit-norm text embedding spaces under
cosine-threshold retrieval}. We do not claim the law governs
biological memory, multimodal embeddings, non-contrastive encoders,
or retrieval with learned scoring heads. The cognitive-science
literature motivated our framing
(\S\ref{sec:related:cls}) and the theorem offers a candidate
geometric reading of the CLS picture, but we do not present it as a
model of biological memory. Transfer to those settings is a
research project, not a corollary of our result.

Within the scope above, the tested surface is concrete: six English
text corpora, six sentence encoders, five consolidation operators,
five learned-quantization baselines, one reader model
(Llama-3.1-70B-Instruct), and corpus sizes from $10$K to $1$M
vectors. The universality claims in \S\ref{sec:template_collapse}
should be read with this tested surface in mind.

\subsection{The central open problem: anisotropic refinement}
\label{sec:disc:open}

The single most consequential weakness of the current theorem is
that its cap-volume argument is \emph{isotropic}: it treats all
directions on $\SphereD$ symmetrically. Real clusters are not
isotropic --- variance concentrates along a few principal
directions, exactly the situation the cluster's $\deff$ quantifies.
This isotropy assumption is what makes our calibrated $c_1$ blow
up to $10^{10}$ in the extreme-coherence corner of the spread
regime (\S\ref{sec:e1:calibration}), and it is why we read
spread-regime evidence throughout the paper as
\emph{observational} support for the law's qualitative predictions
rather than as a test of its quantitative form.

We believe an anisotropic refinement --- replacing the cap volume
with an ellipsoidal volume whose axes are the cluster's principal
directions --- should tighten the bound dramatically in the spread
regime. Constructing that bound is the paper's primary theoretical
open problem. We have not done it; we explicitly raise it so that a
future reader can. The shape-of-the-refinement question is concrete
enough to state as a theorem target: prove a version of
Theorem~\ref{thm:duality} in which $(\thetap/\dbar)^{\deff/2}$ is
replaced by a quantity involving the full eigenvalue profile of
$\Sigma_k$, and show that the calibrated $c_1$ in the spread regime
remains bounded by a polynomial of $\deff$.

\subsection{Other open problems}
\label{sec:disc:other-open}

\paragraph{Streaming consolidation.} All our consolidation operators
run on a fixed corpus. The production version of consolidation is
streaming. An online variant of the centroid operator with amortised
cost and a regret bound against the offline optimum is a natural
next step; the law we prove suggests the regret structure should
track the cluster-level $\dbar_k$.

\paragraph{Multi-level hierarchical consolidation.} E7 shows that
the right centroid behaviour depends on cluster granularity. A
hierarchical index that consolidates at multiple levels
simultaneously --- paragraph centroids within document medoids
within topic clusters --- could unify the ranking. We have not
built this.

\paragraph{Billion-scale empirical validation.} E5 reaches $1$M;
production is $10^8$--$10^9$. The cluster-then-consolidate pipeline
scales linearly with $n$ in cluster time and roughly quadratically
in query time. The scaling slopes we observe through $1$M are
consistent with smooth extrapolation, but extrapolation is not
measurement.

\paragraph{Why this matters.} If the law and centroid lesson
transfer beyond our tested surface, a billion-scale RAG index can be
compressed by $5$--$10$\texttimes{} with the fixed-centroid
operator while preserving identity, and the savings target the
dominant cost of deployed retrieval systems. If they do not
transfer --- if the regime on production corpora is different from
the one we measured --- the same law gives a cheap diagnostic that
surfaces the mismatch before deployment. Either outcome is
informative; the law converts a tuning exercise into a
measurement.

\section{Limitations}
\label{sec:limits}

We flag six specific limitations a reader should hold in mind when
applying the results.

\begin{enumerate}[leftmargin=1.8em, itemsep=4pt]
\item \textbf{L1. The theorem is vacuous in the extreme-coherence
spread regime.} The $95\%$ calibrated $c_1$ is $\approx 0.05$ in tight cells
and $\sim 10^{10}$ in the high-$\dbar$, low-$\deff$ corner. The
spread-regime evidence throughout the paper should be read as
observational support for the theorem's qualitative predictions
(ordering, monotonicity), not as a test of its quantitative form.
We consider the anisotropic refinement the paper's primary open
theoretical problem (\S\ref{sec:disc:open}), and we foreground it
rather than bury it.

\item \textbf{L2. Tested scope is text-centric.} All experiments use
contrastive English-text encoders and cosine-threshold retrieval.
We do not claim the law applies to multimodal embeddings, image
retrieval, speech embeddings, non-contrastive encoders, or
retrieval pipelines that re-rank with a learned scoring head. We
expect the shape of the law to transfer because the cap-volume
argument is geometric, but we have not measured it.

\item \textbf{L3. Compute-gated cells in E7.} We report $10$ of a
possible $15$ downstream QA cells. The missing cells
(selective-prune on NQ, GAC and centroid on HotpotQA,
no-consolidation and medoid on PopQA) are omitted because of
70B inference budget. Their omission weakens any attempt to fit a
full two-parameter interaction model; we treat E7 as directional
evidence rather than a fully-crossed factorial.

\item \textbf{L4. Seed count at $1$M in E5.} The three $1$M cells
were run at one seed each. We state their accuracies as point
evidence that $10$K $\to$ $100$K trends continue, not as
independent measurements with confidence intervals.

\item \textbf{L5. E9 temporal standard deviations.} Earlier drafts
summarised MRR@$20$ across-epoch variability as ``$<0.01$
everywhere''. The correct statement is that the \emph{median}
across-epoch s.d.\ is $0.028$ and $98$ of $105$ cells have s.d.\
$<0.05$, with the highest single cell at $0.108$. The strategy
ranking is nonetheless invariant: centroid is top in every
(corpus, encoder, epoch) combination we measured.

\item \textbf{L6. E4 paraphrase procedure.} The paraphrase
perturbation is embedding-space isotropic Gaussian noise calibrated
to the empirical paraphrase cosine ($\approx 0.92$). Natural
text-level paraphrase may behave differently, though on the two
corpora where we can check (MS MARCO, HotpotQA) the difference is
$< 0.03$ identity points.
\end{enumerate}

\section{Data and code availability}
\label{sec:data}

All code, raw shard JSONL outputs, aggregated parquet tables, and
figure-generation scripts are available at
\url{https://github.com/niashwin/geometry-of-consolidation}. The
repository pins each experiment's provenance: E1 (16{,}000 cells,
synthetic grid), E2 (503 cells, 6 corpora), E3 (138 of 150 cells,
learned-quantization Pareto), E4 (90 cells, paraphrase), E5 (23 cells,
scale with 3 at $1$M), E6 (126 cells, 7-router ablation), E7 (10
cells, Llama-70B downstream; 5 cells, 8B reference), E8 (198 cells,
6-encoder universality), E9 (525 cells, 5-epoch temporal). The total
is $17{,}813$ cells plus $5$ reference cells, all reproducible from
the included shards.
